% Source: https://artofproblemsolving.com/community/c3h1481175_little_help
% Stable identity: aops-topic-1481175-q01
The flag has area $4\cdot2=8$ square feet. The two white rectangles have area $2\cdot(1\cdot2)=4$, and the symbol covers $\dfrac{1}{3}\cdot4=\dfrac{4}{3}$ of the black square. So the white area is $4+\dfrac{4}{3}=\dfrac{16}{3}$, and the white fraction is
\[
\dfrac{16}{3}\div8=\boxed{\dfrac{2}{3}} .
\]
(The symbol's shape does not matter; in the sketch it is drawn as a white disc.)

\begin{asy}
usepackage("amsmath");
size(5cm);
defaultpen(fontsize(10pt));
fill((1,0)--(3,0)--(3,2)--(1,2)--cycle,black);
real rr=sqrt(4/(3*pi)); fill(circle((2,1),rr),white);
draw((0,0)--(4,0)--(4,2)--(0,2)--cycle,linewidth(0.9)); draw((1,0)--(1,2)); draw((3,0)--(3,2));
label("$1$",(0.5,2),N); label("$2$",(2,2),N); label("$1$",(3.5,2),N); label("$2$",(4,1),E);
label("symbol",(2,1),fontsize(8pt));
label("\mbox{}",(min(currentpicture,user=true)+max(currentpicture,user=true))/2,black); // keeps the text after the figure black when LaTeX typesets the labels inline
\end{asy}
